\documentclass[11pt]{article}
\usepackage[margin=1in]{geometry}
\usepackage{amsmath,amssymb,amsthm}
\usepackage{graphicx}
\usepackage{hyperref}
\usepackage{enumitem}

\hypersetup{
  colorlinks=true,
  linkcolor=blue,
  urlcolor=blue,
  citecolor=blue
}
\emergencystretch=2em

\newtheorem{theorem}{Theorem}
\newtheorem{remark}{Remark}

\newcommand{\ellone}{\ell_1}
\newcommand{\ellpq}{\ell_{p,q}}
\newcommand{\chimon}{\chi_{\mathrm{mon}}}
\newcommand{\Pm}{\mathcal P_m}
\newcommand{\id}{\mathrm{id}}

\title{A Lorentz-Scale Partial Result for Lempert's Problem}
\author{Automatic Math Discovery Smoke Test}
\date{June 21, 2026}

\begin{document}
\maketitle

\section*{Source and Target}

The source paper is Andreas Defant, Daniel Galicer, Martin Mansilla,
Mieczyslaw Mastylo, and Santiago Muro,
\emph{Local constants and Bohr's phenomenon for Banach spaces of analytic
polynomials}, arXiv:2509.26326.

After Theorem 2.21, on page 28 of the arXiv PDF, the authors recall Lempert's
problem: whether the projection constants in Theorem 2.21 can be replaced by
the unconditional basis constants of the monomial basis.  More concretely, for
a Banach sequence lattice \(X\), the open question is whether
\[
  X\equiv \ellone
  \quad\Longleftrightarrow\quad
  \sup_{m,n}\chimon(\Pm(X_n))^{1/m}<\infty.
\]
They also recall a theorem of Bayart which gives the following necessary
condition for the right-hand side: for every \(\varepsilon>0\), there is
\(D(\varepsilon)>0\) such that
\[
  \|\id:X_n\to \ell_1^n\|
  \le D(\varepsilon)(\log\log n)^\varepsilon
  \qquad(n\ge 3).
\]

\begin{center}
\includegraphics[width=\textwidth]{figures/open_problem_crop.png}
\end{center}

\section*{Result}

We settle this question on the classical Lorentz scale.

\begin{theorem}
Let \(X=\ell_{p,q}\), where \(1\le p<\infty\) and
\(1\le q\le\infty\), with the usual Lorentz sequence norm.  If
\[
  \sup_{m,n}\chimon(\Pm(X_n))^{1/m}<\infty,
\]
then \(p=q=1\).  Hence, among Lorentz sequence spaces, Lempert's condition can
hold only for \(X=\ellone\).  Conversely, \(X=\ell_{1,1}=\ellone\) is the
known positive model in Lempert's problem.
\end{theorem}

\section*{Proof Intuition}

Bayart's theorem says that any counterexample to Lempert's expected
classification must look almost like \(\ellone\) on finite sections: the
identity \(X_n\to\ell_1^n\) may grow, but no faster than every tiny power of
\(\log\log n\).  Lorentz spaces miss this threshold in a very visible way.
For \(p>1\), the all-ones vector already gives polynomial growth
\(n^{1/p'}\).  For \(p=1\) but \(q>1\), the harmonic vector
\((1,1/2,\ldots,1/n)\) gives a positive power of \(\log n\).  These lower
bounds are far too large to be compatible with Bayart's necessary condition.

\section*{Proof}

Assume
\[
  \sup_{m,n}\chimon(\Pm(X_n))^{1/m}<\infty.
\]
By the Bayart theorem recalled in the source paper, for every
\(\varepsilon>0\) there is \(D(\varepsilon)>0\) such that
\begin{equation}\label{eq:bayart}
  \|\id:X_n\to \ell_1^n\|
  \le D(\varepsilon)(\log\log n)^\varepsilon
  \qquad(n\ge 3).
\end{equation}

We show that \eqref{eq:bayart} fails for every Lorentz space except
\(\ell_{1,1}\).

First suppose \(p>1\).  Let
\[
  u^{(n)}=(1,\ldots,1,0,\ldots)
\]
with \(n\) nonzero coordinates.  The Lorentz fundamental function satisfies
\[
  \|u^{(n)}\|_{\ell_{p,q}}\asymp n^{1/p},
\]
with constants depending only on \(p,q\), also when \(q=\infty\).  Since
\(\|u^{(n)}\|_1=n\), it follows that
\[
  \|\id:\ell_{p,q}^n\to\ell_1^n\|
  \ge c_{p,q} n^{1-1/p}
  = c_{p,q} n^{1/p'}.
\]
This contradicts \eqref{eq:bayart} for any fixed
\(\varepsilon>0\), once \(n\) is large.

It remains to consider \(p=1\).  If \(q=1\), then
\(\ell_{1,1}=\ellone\).  Let \(1<q<\infty\), and set
\[
  h^{(n)}=(1,1/2,\ldots,1/n,0,\ldots).
\]
For the usual Lorentz norm,
\[
  \|h^{(n)}\|_{\ell_{1,q}}
  =
  \left(\sum_{k=1}^n (1/k)^q k^{q-1}\right)^{1/q}
  =
  \left(\sum_{k=1}^n \frac1k\right)^{1/q}
  \asymp (\log n)^{1/q}.
\]
At the same time,
\[
  \|h^{(n)}\|_1=\sum_{k=1}^n\frac1k\asymp \log n.
\]
Thus
\[
  \|\id:\ell_{1,q}^n\to\ell_1^n\|
  \ge c_q(\log n)^{1-1/q}
  =c_q(\log n)^{1/q'}.
\]
Choosing, for instance, \(0<\varepsilon<1/q'\), this again contradicts
\eqref{eq:bayart} for large \(n\).

Finally, if \(p=1\) and \(q=\infty\), the same vector \(h^{(n)}\) satisfies
\[
  \|h^{(n)}\|_{\ell_{1,\infty}}
  =\sup_{1\le k\le n} k\cdot \frac1k=1,
\qquad
  \|h^{(n)}\|_1\asymp \log n.
\]
Therefore
\[
  \|\id:\ell_{1,\infty}^n\to\ell_1^n\|\ge c\log n,
\]
contradicting \eqref{eq:bayart}.

Hence the bounded-root monomial unconditionality condition cannot hold for
any Lorentz sequence space \(\ell_{p,q}\) except \(\ell_{1,1}=\ellone\), as
claimed.

\section*{Verification Notes}

The proof uses only two inputs beyond the source definitions:
\begin{itemize}[leftmargin=*]
\item Bayart's necessary condition, quoted in the source paper after Theorem
2.21 and attributed there to \emph{Maximum modulus of random polynomials}.
\item The elementary Lorentz norm computations
\(\|(1,\ldots,1)\|_{\ell_{p,q}^n}\asymp n^{1/p}\) and
\(\|(1,1/2,\ldots,1/n)\|_{\ell_{1,q}^n}\asymp(\log n)^{1/q}\) for
\(1<q<\infty\), with the weak-\(\ell_1\) norm equal to \(1\) in the case
\(q=\infty\).
\end{itemize}

No numerical computation is used.

\section*{Scope and Human Review Recommendation}

This is a partial result only.  It does not solve Lempert's problem outside
the Lorentz family, and it does not improve Bayart's general necessary
condition.  It records that the Lorentz spaces naturally present in the source
paper cannot supply a counterexample to Lempert's expected classification.

Human review should check the exact formulation of the Lorentz norm convention
used for \(q=\infty\), the dependence of the constants on \(p,q\), and whether
the source's word ``\(X\equiv\ell_1\)'' is intended as isometric equality of
the sequence lattice or Banach-space isomorphism.  Under the standard Lorentz
normalizations, the finite-section identity lower bounds are invariant up to
absolute constants and the conclusion is unchanged.

\section*{Search Notes}

Run indexes were searched for arXiv:2509.26326, Lempert, monomial basis,
Bayart, \(\log\log\), and Lorentz variants.  External searches on
2026-06-21 for ``Lempert's problem monomial basis'', ``Lempert's problem log
log Bayart'', and ``Lempert Lorentz monomial'' found no later full solution or
duplicate Lorentz-scale packet.

\section*{References}

\begin{itemize}[leftmargin=*]
\item A. Defant, D. Galicer, M. Mansilla, M. Mastylo, and S. Muro,
\emph{Local constants and Bohr's phenomenon for Banach spaces of analytic
polynomials}, arXiv:2509.26326.
\item F. Bayart, \emph{Maximum modulus of random polynomials}, Q. J. Math.
63 (2012), no. 1, 21--39.
\item J. Lindenstrauss and L. Tzafriri, \emph{Classical Banach Spaces I:
Sequence Spaces}, Springer, 1977, for standard Lorentz sequence-space
notation and finite-section estimates.
\end{itemize}

\end{document}
